% answers_short.tex —— 2014 年·陕西卷·文科 · 简短答案速查
% 由 tools/make_short_answers.py 自 answers.tex 的 \gksolhead/\gkanswer 提取，
% 并与 2014/2014_数学真题_陕西文_答案要点.md 逐题比对后生成，勿手改（改源 MD 后重跑）。

\begin{gknotice}
\textbf{使用说明}：本节只给最终答案便于快速核对，完整推导见解析版。答案为编写者独立推导，\textbf{非官方标准答案}。
\end{gknotice}

\gkpart{选择题}
\begin{center}\begin{tabular}{ccccccccccc}
\toprule
题号 & 1 & 2 & 3 & 4 & 5 & 6 & 7 & 8 & 9 & 10 \\
\midrule
答案 & D & B & A & C & C & C & D & A & D & A \\
\bottomrule
\end{tabular}\end{center}

\gkpart{填空题}
\begin{itemize}[leftmargin=2.6em]
  \item[\textbf{11.}] $x=-1$。抛物线 $y^2=2px$ 的准线为 $x=-\dfrac p2$，此处 $2p=4$。
  \item[\textbf{12.}] $\sqrt{10}$。$4^a=2$ 得 $a=\dfrac12$，故 $x=10^{1/2}=\sqrt{10}$。
  \item[\textbf{13.}] $\dfrac12$。由 $\boldsymbol a\cdot\boldsymbol b=0$ 得 $\sin2\theta=\cos^2\theta$，即 $2\sin\theta\cos\theta=\cos^2\theta$，故 $\tan\theta=\dfrac12$。
  \item[\textbf{14.}] $\dfrac{x}{1+2014x}$。由 $f_n(x)=\dfrac{x}{1+nx}$ 归纳可得。
  \item[\textbf{15.}] A 为 $\sqrt5$；B 为 $3$；C 为 $1$。A 用柯西不等式；B 由 $\triangle AEF\sim\triangle ACB$ 且相似比为 $\dfrac12$；C 化极坐标方程为 $-x+\sqrt3y=2$ 后求点到直线距离。
\end{itemize}

\gkpart{解答题}
\begin{itemize}[leftmargin=2.6em]
  \item[\textbf{16.}] (1) 见解析版；(2) $\dfrac34$。由 $b^2=ac=2a^2$、$c=2a$，得 $\cos B=\dfrac{a^2+4a^2-2a^2}{4a^2}=\dfrac34$
  \item[\textbf{17.}] (1) $\dfrac23$。由三视图得 $DA=1$，$DB=DC=2$ 且两两垂直，$V=\dfrac13\cdot\dfrac12\cdot2\cdot2\cdot1=\dfrac23$；(2) 见解析版
  \item[\textbf{18.}] (1) $2\sqrt2$。$\overrightarrow{OP}=\dfrac23(\overrightarrow{AB}+\overrightarrow{AC})=(2,2)$，故 $|\overrightarrow{OP}|=2\sqrt2$；(2) $m-n=y-x$，最大值为 $1$。由 $x=m+2n$，$y=2m+n$ 解出 $m-n=y-x$，在三角形三个顶点处比较得最大值 $1$
  \item[\textbf{19.}] (1) $0.27$。赔付金额大于 $2800$ 元即 $3000$ 元或 $4000$ 元，频率为 $\dfrac{150+120}{1000}=0.27$；(2) $0.024$。所求为“新司机且获赔 $4000$ 元”的概率，$\dfrac{120}{1000}\times20\%=0.024$
  \item[\textbf{20.}] (1) $\dfrac{x^2}{4}+\dfrac{y^2}{3}=1$。由过点 $(0,\sqrt3)$ 得 $b=\sqrt3$，由 $e=\dfrac12$ 得 $a=2$；(2) $y=-\dfrac12x+\dfrac{\sqrt3}{3}$ 或 $y=-\dfrac12x-\dfrac{\sqrt3}{3}$。由弦长比解得 $m^2=\dfrac13$
  \item[\textbf{21.}] (1) $2$。$f'(x)=\dfrac{x-\mathrm e}{x^2}$，极小值点为 $x=\mathrm e$，$f(\mathrm e)=1+1=2$；(2) 当 $m>\dfrac23$ 时零点个数为 $0$；当 $m=\dfrac23$ 或 $m\leqslant0$ 时为 $1$；当 $0<m<\dfrac23$ 时为 $2$。零点个数即 $h(x)=x^3-3x+3m$ 在 $(0,+\infty)$ 上的零点个数；(3) $\left[\dfrac14,+\infty\right)$。由差商极限得 $f'(x)\leqslant1$ 对一切 $x>0$ 成立，即 $m\geqslant x-x^2$，故 $m\geqslant\dfrac14$
\end{itemize}
